% Generated draft: requires mathematical and source review.
\documentclass[11pt]{article}
\usepackage[margin=25mm]{geometry}
\usepackage{fontspec}
\setmainfont{Latin Modern Roman}
\usepackage{amsmath,amssymb,mathtools,unicode-math}
\setmathfont{Latin Modern Math}
\usepackage{enumitem,needspace,xurl,hyperref}
\setlength{\parindent}{0pt}
\setlength{\parskip}{5pt}
\setlength{\emergencystretch}{3em}
\newcommand{\N}{\mathbb{N}}
\newcommand{\Z}{\mathbb{Z}}
\newcommand{\Q}{\mathbb{Q}}
\newcommand{\R}{\mathbb{R}}
\newcommand{\C}{\mathbb{C}}
\newcommand{\F}{\mathbb{F}}
\newcommand{\A}{\mathbb{A}}
\newcommand{\PP}{\mathbb{P}}
\newcommand{\E}{\mathbb{E}}
\newcommand{\eps}{\varepsilon}
\newcommand{\dd}{\,\mathrm{d}}
\begin{document}
\section*{UnsolvedMath ID: 20002021}
\textbf{Source problem code:} catalogue entry
\subsection{Definitions and mathematical statement}
Michael Eastwood's Problem 3 asks, without further qualification, whether there is a global ambient metric construction. The paired aid distinguishes: construction formal along the entire conformal base, an exact smooth Ricci-flat ambient metric on a domain, and a geodesically complete global extension. These are different assertions.

In the usual conformal dimension \(n\geq3\), an ambient normal form for a representative \(g\) is
\[
 \widetilde g=2\rho\,dt^2+2t\,dt\,d\rho+t^2g_\rho,
 \qquad t>0,\quad g_0=g.
\]
It is homogeneous of degree two under scaling of \(t\) and restricts to the tautological metric on the cone at \(\rho=0\). The Fefferman--Graham formal construction is global along the base in its appropriate formal sense. In even dimensions its obstruction and allowed order/ambiguity must be retained; a nonzero ambient obstruction prevents an exact smooth Ricci-flat construction of the desired type. Formal base-globality is therefore not unconditional existence of an exact smooth all-order metric or a complete manifold.

For an Einstein representative with
\[
 \operatorname{Ric}(g)=2\lambda(n-1)g,
\]
the known exact family reported and verified in the aid is
\[
 g_\rho=(1+\lambda\rho)^2g,\qquad
 \widetilde g=2\rho\,dt^2+2t\,dt\,d\rho
                  +t^2(1+\lambda\rho)^2g.
\]
It is Ricci flat on \((0,\infty)\times M\times I_\lambda\), where the maximal connected nondegenerate normal-form interval containing \(0\) is
\[
 I_\lambda=\begin{cases}
 \mathbb R,&\lambda=0,\\
 (-1/\lambda,\infty),&\lambda>0,\\
 (-\infty,-1/\lambda),&\lambda<0.
 \end{cases}
\]
This construction is exact over the entire base for this special class. Maximality here is only within the displayed connected nondegenerate normal-form domain, not maximality among all coordinate-independent ambient extensions.

The aid supplies an incompleteness diagnostic on that domain. With \(x,\rho\) fixed, coordinate \(t\)-lines are affine geodesics and reach the missing endpoint \(t=0\) in finite affine parameter. For \(\lambda\ne0\), coordinate \(\rho\)-lines with \(t,x\) fixed are also affine geodesics and reach \(\rho=-1/\lambda\) in finite parameter. Thus this exact base-global normal-form domain is geodesically incomplete, including the \(\lambda=0\) case via \(t\)-lines. The conclusion does not prohibit a further extension in other coordinates, nor establish a general no-complete-ambient theorem. The original word global must be specified before these special-case results can be regarded as an answer.
\subsection{Short English statement}
The supplied aid separates formal base-global construction, exact smooth construction and completeness. An explicit Einstein ambient metric is Ricci flat on a maximal nondegenerate normal-form interval, yet that domain is incomplete. Even-dimensional obstruction limits general exact smooth constructions; a complete extension is a separate question.
\subsection{Sources}
\begin{itemize}
\item Local source record: catalogue entry, catalogue entry; canonical \texttt{aim-geometry-notes.json}, \texttt{notes[358]}, Problem 3, under Michael Eastwood's contribution. Title: ``Three meanings of a global ambient metric.''
\item AIM workshop \emph{Conformal structure in geometry, analysis, and physics}, supplied reference \url{https://aimath.org/WWN/confstruct/confstruct.pdf}. The original question is exactly ``Is there a global ambient metric construction?'' Original remarks and literature fields are empty; this draft does not access a PDF.
\item Supplied research summary invokes the Fefferman--Graham formal construction and ambient obstruction, and a known Einstein normal-form family. No full bibliographic entries or theorem numbers are supplied. Attempt 1 supplies the interval and affine-coordinate-geodesic diagnostic.
\item Registry dated 14 August 2026: ``partially\_solved,'' ``PARTIAL-PROGRESS,'' statement status ``exact''; attempt partial, rigor ``proved,'' novelty confidence low. The machine-generated, non-peer-reviewed aid requires independent verification before citation.
\end{itemize}
\end{document}
